%!TEX root = ../exa-ma-deliverables.tex
\chapter{Project Management and Activity Report (D0.2)}
\label{chap:management}

\section{Scope and annual reporting periods}
\label{sec:management-scope}

This chapter provides the reporting structure for Deliverable D0.2,
\emph{Project Management and Activity Report}, within the Exa-MA
mega-deliverable. It brings together the annual assessment of project
management, staffing, collective activities and external partnerships before
the detailed scientific reports of WP1--WP6 and the WP7 activities.

The Exa-MA scientific project document (Section 3.3, WP0) assigns project
management to WP0, led by UNISTRA. D0.2 covers coordination and administrative
management (T0.1), scientific and technical management (T0.2), provision of
computing resources (T0.3), and communication (T0.4), with annual deliveries
planned at M12, M24, M36, M48 and M60.
% Source: PC1-ExaMA - Document Scientifique.md, Section 3.3, WP0 task table.
% This source describes the planned mandate, not evidence of completed work.

\textbf{Draft status.} The annual sections below are editorial placeholders;
they do not establish that an activity took place or that a partnership was
concluded. The coordination team must confirm the calendar dates and reporting
cut-off for each project-year period, including any approved schedule changes.
Missing information must be distinguished from a confirmed absence of activity.

\section{Common reporting framework}
\label{sec:management-framework}

Each annual assessment should explain what happened, how it affected project
delivery, which decisions were taken and what remains to be done. Entries should
identify the relevant period, partner or WP, responsible role, status and dated
supporting evidence. Planned actions and completed activities should be reported
separately.

\subsection{Governance, coordination and resources}
Summarize governance meetings, participation, major decisions and their follow-up;
coordination across WPs and with NumPEx; deliverable reviews and milestones;
administrative and financial resource monitoring; and changes to the work plan.
Describe difficulties, risks, corrective actions and their effects on the project.
Support the assessment with meeting minutes, decision records and validated
resource summaries.

\subsection{Human resources: recruitment, arrivals and departures}
Report staffing by partner, WP and role (doctoral researchers, postdoctoral
researchers, engineers, interns and permanent staff). Record opening and closing
headcounts, arrivals, departures, internal transfers, vacancies and recruitment
delays, with dates or periods. Where available, distinguish headcounts from
full-time equivalents and person-months. Explain onboarding, supervision,
handover and continuity measures, and the impact of staffing changes on delivery.
Use consolidated information appropriate for this public report.

\subsection{Events, communication and collective activities}
List project meetings, workshops, seminars, schools, conferences and outreach
activities organized, co-organized or attended. Record dates, locations, the
project's role, audience, participation where documented, outputs and follow-up.
Assess their contribution to coordination, dissemination and community building.
For training events, refer to \Cref{chap:training} for educational content and
reuse the same event identifiers to avoid double-counting.

\subsection{Industrial relations and partnerships}
Identify the companies involved, the purpose of each interaction, the Exa-MA
contact role and the relevant WPs. Distinguish exploratory contacts, discussions,
formalized agreements and active collaborations. Describe joint use cases,
co-development, technology transfer or joint events only when supported by
evidence; record outcomes, obstacles and next steps without treating an initial
contact as an established partnership.

\subsection{Research institutions and computing centres}
Describe cooperation with external universities, major research institutions,
national and European computing centres, and related projects. Distinguish
external cooperation from the consortium's internal coordination. Record joint
research, visits, shared platforms, computing allocations and technical support
where documented. For computing resources, distinguish requests, awards and
actual use, and explain their contribution to the work plan. Cross-reference
the relevant scientific, application or benchmark sections for technical results.

\subsection{Annual assessment and next-period priorities}
Conclude each annual report with achievements, unresolved issues, lessons learned
and actions for the next period. Give each action an owner, target date and
expected outcome; carry unresolved actions into the following year's assessment.

% Annual records are authored independently to preserve their reporting history.
% Replace each placeholder with a validated account and dated evidence.
%!TEX root = ../../exa-ma-d7.1.tex
\clearpage
\section{Year 1: M1--M12}
\label{sec:management-year-1}

\textbf{Reporting period and validation.}
\emph{To complete: calendar start and end dates, reporting cut-off, contributors,
coordination reviewer and validation date.}

\subsection{Governance, coordination and resources}
\emph{To complete: annual management assessment, meetings and decisions,
milestone and deliverable follow-up, resource situation, difficulties and
corrective actions, with dated evidence.}

\subsection{Human resources: arrivals and departures}
\emph{To complete: opening and closing staffing by partner/WP/role, recruitment,
arrivals, departures, transfers and vacancies; effects on delivery and handover
measures, with a validated staffing summary.}

\subsection{Events and communication}
\emph{To complete: dated event inventory, organizers, Exa-MA role, audiences,
documented participation, outputs and follow-up; links to programmes, materials
and the training chapter where applicable.}

\subsection{Industrial relations and partnerships}
\emph{To complete: companies, purpose, contact roles, collaboration status,
activities and outcomes during this period; dated supporting evidence and
next steps.}

\subsection{Research institutions and computing centres}
\emph{To complete: external institutions and centres, joint activities,
resource requests/awards/use, results for the project and evidence;
distinguish consortium coordination from external partnerships.}

\subsection{Lessons learned and next-period actions}
\emph{To complete: overall assessment, follow-up of earlier actions where
applicable, outstanding risks, lessons learned and next-period priorities,
each with an owner, target date and expected outcome.}

%!TEX root = ../../exa-ma-d7.1.tex
\clearpage
\section{Year 2: M13--M24}
\label{sec:management-year-2}

\textbf{Reporting period and validation.}
\emph{To complete: calendar start and end dates, reporting cut-off, contributors,
coordination reviewer and validation date.}

\subsection{Governance, coordination and resources}
\emph{To complete: annual management assessment, meetings and decisions,
milestone and deliverable follow-up, resource situation, difficulties and
corrective actions, with dated evidence.}

\subsection{Human resources: arrivals and departures}
\emph{To complete: opening and closing staffing by partner/WP/role, recruitment,
arrivals, departures, transfers and vacancies; effects on delivery and handover
measures, with a validated staffing summary.}

\subsection{Events and communication}
\emph{To complete: dated event inventory, organizers, Exa-MA role, audiences,
documented participation, outputs and follow-up; links to programmes, materials
and the training chapter where applicable.}

\subsection{Industrial relations and partnerships}
\emph{To complete: companies, purpose, contact roles, collaboration status,
activities and outcomes during this period; dated supporting evidence and
next steps.}

\subsection{Research institutions and computing centres}
\emph{To complete: external institutions and centres, joint activities,
resource requests/awards/use, results for the project and evidence;
distinguish consortium coordination from external partnerships.}

\subsection{Lessons learned and next-period actions}
\emph{To complete: overall assessment, follow-up of earlier actions where
applicable, outstanding risks, lessons learned and next-period priorities,
each with an owner, target date and expected outcome.}

%!TEX root = ../../exa-ma-deliverables.tex
\clearpage
\section{Year 3: M25--M36}
\label{sec:management-year-3}

\textbf{Reporting period and validation.}
\emph{To complete: calendar start and end dates, reporting cut-off, contributors,
coordination reviewer and validation date.}

\subsection{Governance, coordination and resources}
\emph{To complete: annual management assessment, meetings and decisions,
milestone and deliverable follow-up, resource situation, difficulties and
corrective actions, with dated evidence.}

\subsection{Human resources: arrivals and departures}
\emph{To complete: opening and closing staffing by partner/WP/role, recruitment,
arrivals, departures, transfers and vacancies; effects on delivery and handover
measures, with a validated staffing summary.}

\subsection{Events and communication}
\emph{To complete: dated event inventory, organizers, Exa-MA role, audiences,
documented participation, outputs and follow-up; links to programmes, materials
and the training chapter where applicable.}

\subsection{Industrial relations and partnerships}
\emph{To complete: companies, purpose, contact roles, collaboration status,
activities and outcomes during this period; dated supporting evidence and
next steps.}

\subsection{Research institutions and computing centres}
\emph{To complete: external institutions and centres, joint activities,
resource requests/awards/use, results for the project and evidence;
distinguish consortium coordination from external partnerships.}

\subsection{Lessons learned and next-period actions}
\emph{To complete: overall assessment, follow-up of earlier actions where
applicable, outstanding risks, lessons learned and next-period priorities,
each with an owner, target date and expected outcome.}

%!TEX root = ../../exa-ma-d7.1.tex
\clearpage
\section{Year 4: M37--M48}
\label{sec:management-year-4}

\textbf{Reporting period and validation.}
\emph{To complete: calendar start and end dates, reporting cut-off, contributors,
coordination reviewer and validation date.}

\subsection{Governance, coordination and resources}
\emph{To complete: annual management assessment, meetings and decisions,
milestone and deliverable follow-up, resource situation, difficulties and
corrective actions, with dated evidence.}

\subsection{Human resources: arrivals and departures}
\emph{To complete: opening and closing staffing by partner/WP/role, recruitment,
arrivals, departures, transfers and vacancies; effects on delivery and handover
measures, with a validated staffing summary.}

\subsection{Events and communication}
\emph{To complete: dated event inventory, organizers, Exa-MA role, audiences,
documented participation, outputs and follow-up; links to programmes, materials
and the training chapter where applicable.}

\subsection{Industrial relations and partnerships}
\emph{To complete: companies, purpose, contact roles, collaboration status,
activities and outcomes during this period; dated supporting evidence and
next steps.}

\subsection{Research institutions and computing centres}
\emph{To complete: external institutions and centres, joint activities,
resource requests/awards/use, results for the project and evidence;
distinguish consortium coordination from external partnerships.}

\subsection{Lessons learned and next-period actions}
\emph{To complete: overall assessment, follow-up of earlier actions where
applicable, outstanding risks, lessons learned and next-period priorities,
each with an owner, target date and expected outcome.}

%!TEX root = ../../exa-ma-d7.1.tex
\clearpage
\section{Year 5: M49--M60}
\label{sec:management-year-5}

\textbf{Reporting period and validation.}
\emph{To complete: calendar start and end dates, reporting cut-off, contributors,
coordination reviewer and validation date.}

\subsection{Governance, coordination and resources}
\emph{To complete: annual management assessment, meetings and decisions,
milestone and deliverable follow-up, resource situation, difficulties and
corrective actions, with dated evidence.}

\subsection{Human resources: arrivals and departures}
\emph{To complete: opening and closing staffing by partner/WP/role, recruitment,
arrivals, departures, transfers and vacancies; effects on delivery and handover
measures, with a validated staffing summary.}

\subsection{Events and communication}
\emph{To complete: dated event inventory, organizers, Exa-MA role, audiences,
documented participation, outputs and follow-up; links to programmes, materials
and the training chapter where applicable.}

\subsection{Industrial relations and partnerships}
\emph{To complete: companies, purpose, contact roles, collaboration status,
activities and outcomes during this period; dated supporting evidence and
next steps.}

\subsection{Research institutions and computing centres}
\emph{To complete: external institutions and centres, joint activities,
resource requests/awards/use, results for the project and evidence;
distinguish consortium coordination from external partnerships.}

\subsection{Lessons learned and next-period actions}
\emph{To complete: overall assessment, follow-up of earlier actions where
applicable, outstanding risks, lessons learned and next-period priorities,
each with an owner, target date and expected outcome.}

